\documentclass[11pt]{article}

\usepackage{geometry}
\geometry{left=1.1in, right=1.1in, top=1.0in, bottom=0.95in}

% graphicx package, useful for including eps and pdf graphics
\usepackage{graphicx}
\DeclareGraphicsExtensions{.pdf,.png,.jpg}

% basic packages
\usepackage{color}
\usepackage{parskip}
\setlength{\parskip}{0.16cm}
\usepackage{float}
\usepackage{todonotes}
\usepackage{enumitem}
\usepackage{microtype}

\definecolor{green}{rgb}{0.20,0.50,0.48}
\usepackage[hidelinks]{hyperref}
\hypersetup{colorlinks=true,linkcolor=black,citecolor=black,urlcolor=green}

\setlength{\parindent}{0pt} % Remove paragraph indent

\definecolor{brown}{rgb}{0.700,0.150,0.150}
\def\mfc#1{\textcolor{brown}{[#1]}}

\begin{document}

\thispagestyle{empty} % Remove page headers/footers

\mbox{}\hfill
\begin{tabular}{l @{}}
	\includegraphics[width=6.5cm]{fhcc_logo} \\
	Vaccine and Infectious Disease Division \\
	Fred Hutchinson Cancer Center \\
	1100 Fairview Ave N \\
	Seattle, WA 98109, USA \\
	Email: \href{mailto:jhuddles@fredhutch.org}{jhuddles@fredhutch.org} \\
\end{tabular}

\vspace{0.1in} % Vertical skip between sender/receiver address

\begin{tabular}{@{} l}
  \today
\end{tabular}

\vspace{0.1in} % Vertical skip between receive address and letter opening

Dear Editors,

\medskip % Vertical skip between letter opening and letter body

Please find attached our manuscript entitled ``Dimensionality reduction distills complex evolutionary relationships in seasonal influenza and SARS-CoV-2''.
We would be grateful if you considered it for publication in \emph{Virus Evolution}.

Phylogenetic approaches are now common tools for public health practitioners who seek to understand transmission dynamics and viral evolution during outbreaks, epidemics, and pandemics.
However, viruses that cause substantial human disease like seasonal influenza and SARS-CoV-2 evolve through processes of reassortment and recombination that violate the assumption of shared evolutionary history at the core of phylogenetics.
Just as importantly, phylogenies can be ineffective visualizations of genetic relationships for an audience that isn't already familiar with how to interpret them.
For public health responses that rely on measuring and visualizing pairwise genetic distances between viral samples, phylogenetic methods that infer the complete evolutionary history of samples are unnecessarily complex.

We present an alternative approach to phylogenetics where we apply standard dimensionality reduction methods (PCA, MDS, t-SNE, and UMAP) to viral genome sequences and evaluate the ability of the resulting low-dimensional embeddings to represent genetic distances.
After first identifying optimal parameters for each embedding method with simulated populations, we applied these methods to recent seasonal influenza A/H3N2 hemagglutinin and SARS-CoV-2 sequences.
We measured how well the embeddings represented pairwise genetic relationships and how well the clusters that we discovered in the embeddings with an unsupervised learning algorithm (HDBSCAN) corresponded to expert clade definitions.
Finally, we quantified the ability of these embeddings and clusters to capture reassortment events between A/H3N2 hemagglutinin and neuraminidase genes and recombinant lineages in SARS-CoV-2.

We find that these biologically-uninformed statistical methods can accurately represent known genetic relationships for both influenza and SARS-CoV-2 including reassortant and recombinant lineages.
The resulting embeddings produce biologically relevant clusters in a readily interpretable 2- or 3-dimensional visualization.
Given the robust performance of these methods across simulated and natural virus populations, distinct evolutionary time periods, and a wide range of method parameters, researchers and practitioners could immediately apply these methods to other viral pathogens.
Our open source implementation of these methods supports such applications by integrating easily with existing workflows for genomic epidemiology.
We expect this general approach to provide the foundation for future methods that formally identify reassortant or recombinant lineages.

This work builds on a long history publications that have applied dimensionality reduction methods to biological systems and rigorously evaluated the accuracy and interpretability of such methods.
However, this study is the first such detailed analysis of how well these methods capture monophyletic clades, reassortment, and recombination in viral genomes.
Our findings highlight the value of continuing this line of research and should be relevant to an audience of evolutionary biologists, infectious disease researchers, and computational biologists as well as public health practitioners.
For these reasons, we would appreciate your consideration for publication in \emph{Virus Evolution}.

\vspace{0.3in} % Vertical skip between letter closing and signature start

Sincerely, \newline
\vspace{0.05in} \newline
John Huddleston, PhD \newline
Staff Scientist \newline
Vaccine and Infectious Disease Division \newline
Fred Hutchinson Cancer Center

\end{document}
